\chapter{Comparison of Functions}\label{ch:b2:comparison}

Asymptotic analysis --- the art of replacing a complicated quantity
by a simple one plus a controlled error --- was begun in the Year 1
volume with Taylor expansions. This chapter makes it a discipline of
its own: expansions along general scales, the series--integral
comparison with its full asymptotic force, Stirling's formula (proved
completely), and the systematic study of implicitly defined
sequences. These techniques are the daily bread of asymptotic
analysis, and every later chapter that estimates anything ---
series, integrals, probabilities --- eats from this table.

\section{Comparison relations and scales}

\begin{definition}\label{def:b2:comparison:landau}
Near a point $a$ ($a \in \R$ or $\pm\infty$), for functions (or
sequences, with $n \to \infty$): $f = o(g)$, $f = O(g)$, $f \sim g$
as in the Year 1 volume. A \emph{comparison
scale}\index{comparison scale} at $a$ is a family of positive
functions, pairwise comparable, totally ordered by $o(\cdot)$ ---
the standard scale at $+\infty$ being
\[
x^{\alpha} (\ln x)^{\beta}
\qquad (\alpha, \beta \in \R),
\]
ordered lexicographically in $(\alpha, \beta)$, refined when needed
by exponentials $\eu^{\gamma x}$.
\end{definition}

\begin{definition}[Asymptotic expansion]\label{def:b2:comparison:expansion}
$f$ admits the \emph{asymptotic expansion}\index{asymptotic
expansion}
\[
f = c_1 \varphi_1 + c_2\varphi_2 + \dots + c_k \varphi_k +
o(\varphi_k)
\qquad (\varphi_{i+1} = o(\varphi_i) \text{ in the scale})
\]
when the successive remainders satisfy the displayed estimates. The
coefficients are then unique: $c_1 = \lim f/\varphi_1$, and
inductively $c_{i+1} = \lim\,(f - \sum_{j \leq i}
c_j\varphi_j)/\varphi_{i+1}$.
\end{definition}

\begin{example}
Taylor expansions are \omterm{def:b2:comparison:expansion}{asymptotic expansions} along the scale $(x -
a)^k$ at $a$. But the notion is strictly wider: at $+\infty$,
\[
\frac{1}{x - \ln x}
= \frac1x \cdot \frac{1}{1 - \frac{\ln x}{x}}
= \frac1x + \frac{\ln x}{x^2} + o\Bigl(\frac{\ln x}{x^2}\Bigr),
\]
an expansion along the mixed scale --- no Taylor theorem applies,
only the geometric expansion and the calculus of $o$'s.
\end{example}

\begin{example}[The standard scale is really ordered]\label{ex:b2:comparison:scaleorder}
The lexicographic claim of \cref{def:b2:comparison:landau}
needs one line of proof per case. Compare
$x^{\alpha}(\ln x)^{\beta}$ and $x^{\alpha'}(\ln x)^{\beta'}$
at $+\infty$. If $\alpha < \alpha'$: the ratio is
$x^{\alpha - \alpha'}(\ln x)^{\beta - \beta'} \to 0$, because a
negative power of $x$ crushes any power of $\ln x$ (set $x =
\eu^t$: $\eu^{(\alpha - \alpha')t}\,t^{\beta - \beta'} \to 0$
by the exponential-beats-polynomial limit of the Year 1
volume). If $\alpha = \alpha'$ and $\beta < \beta'$: the ratio
is $(\ln x)^{\beta - \beta'} \to 0$ directly. So the pairs
$(\alpha, \beta)$, ordered lexicographically, order the scale
by $o(\cdot)$ --- and the substitution $x = \eu^t$ is the
one-size-fits-all trick for mixed power-log comparisons.
\end{example}

\begin{example}[Ranking a menagerie]\label{ex:b2:comparison:ranking}
Scales must be \emph{ordered}; here is the standard drill. At
$+\infty$, compare $n^{10}$, $\eu^{\sqrt{\ln n}\,\cdot\,\sqrt n}$,
$2^n$ and $n^{\ln n}$ by taking logarithms:
\[
10\ln n
\;\ll\; (\ln n)^2
\;\ll\; \sqrt{n\ln n}
\;\ll\; n\ln 2 ,
\]
where $a_n \ll b_n$ means $a_n = o(b_n)$; the second entry is
$\ln(n^{\ln n})$. Exponentials preserve these strict gaps (if
$\ln u_n - \ln v_n \to -\infty$ then $u_n/v_n \to 0$), so
\[
n^{10} = o\bigl(n^{\ln n}\bigr),
\qquad
n^{\ln n} = o\bigl(\eu^{\sqrt{n\ln n}}\bigr),
\qquad
\eu^{\sqrt{n\ln n}} = o(2^n) .
\]
The moral, twice over: always compare through logarithms
(differences of logs, not ratios of logs), and never conclude
$u_n \sim v_n$ from $\ln u_n \sim \ln v_n$ --- the pair $n^{10}$
and $n^{\ln n}$ has $\ln$-ratio tending to $\infty$, but $2^n$
and $4^n$ have $\ln$-ratio exactly $2$ and are wildly
inequivalent.
\end{example}

\section{Series--integral comparison, asymptotically}

\begin{theorem}\label{thm:b2:comparison:seriesintegral}
Let $f$ be \omterm{def:b2:metric:continuity}{continuous}, positive, \emph{decreasing} on
$\intco{1}{+\infty}$.
\begin{enumerate}
  \item If $\int_1^{\infty} f$ converges, the remainders satisfy
        \[
        \int_{n+1}^{\infty} f \;\leq\; \sum_{k > n} f(k) \;\leq\;
        \int_{n}^{\infty} f .
        \]
  \item If $\int_1^\infty f$ diverges, the partial sums satisfy
        $\sum_{k=1}^{n} f(k) = \int_1^n f + C + o(1)$ for some
        constant $C$: the difference $\sum_{k \leq n} f(k) -
        \int_1^n f$ \emph{converges}.
\end{enumerate}
\end{theorem}

\begin{proof}
The bracketing $f(k+1) \leq \int_k^{k+1} f \leq f(k)$ (decrease) was
the Year 1 device; summing over $k \geq n+1$, resp.\ $k \geq n$,
gives (1). For (2), set $u_k = f(k) - \int_k^{k+1} f$: by the
bracketing, $0 \leq u_k \leq f(k) - f(k+1)$, so the partial sums of
$\sum u_k$ are bounded by the telescoping $f(1) - f(n+1) \leq
f(1)$: the series converges. Moreover the sequence $\bigl(\int_n^{n+1}
f\bigr)_n$ is nonincreasing ($f$ decreases) and nonnegative, hence
convergent. Writing
\[
\sum_{k=1}^{n} f(k) - \int_1^n f
= \sum_{k=1}^{n} u_k + \int_n^{n+1} f ,
\]
the right side converges as $n \to \infty$: the difference converges
to a constant $C$, which is statement (2).
\end{proof}

\begin{example}[The harmonic expansion]\label{ex:b2:comparison:harmonic}
For $f(t) = \frac1t$: $H_n = \ln n + \gamma + o(1)$, recovering
Euler's constant (Year 1 volume) with a cleaner proof. Pushing one
order further (\cref{exo:b2:comparison:3}):
\[
H_n = \ln n + \gamma + \frac{1}{2n} + o\Bigl(\frac1n\Bigr).
\]
Numbers make the gain visible at $n = 10$: $H_{10} =
2.928968\dots$ and $\ln 10 = 2.302585\dots$, so the raw estimate
of $\gamma$ is $H_{10} - \ln 10 = 0.626383$, off by $0.049$;
subtracting the correction $\frac1{20}$ gives $0.576383$,
off from $\gamma = 0.577216$ by only $8.3\cdot10^{-4}$ ---
which is itself the next term $\frac{1}{12\cdot100}$ of the
expansion, as the weekend problem proves (question 8).
\end{example}

\begin{example}[A crude $\ln(n!)$ without Stirling]\label{ex:b2:comparison:lnfactorial}
The bracketing device alone already locates $\ln(n!)$. Since
$\ln$ increases,
\[
\int_{k-1}^{k}\ln t\,\dd t \;\leq\; \ln k \;\leq\;
\int_{k}^{k+1}\ln t\,\dd t ,
\]
and summing over $k = 2, \dots, n$ (with $\int_1^n\ln = n\ln n
- n + 1$):
\[
n\ln n - n + 1 \;\leq\; \ln(n!) \;\leq\; (n+1)\ln(n+1) - n .
\]
Both fences are $n\ln n - n + O(\ln n)$: hence $\ln(n!) = n\ln
n - n + O(\ln n)$, and in particular $\ln(n!) \sim n\ln n$.
What Stirling adds is the next two rungs --- the $\frac12\ln n$
and the constant $\ln\sqrt{2\pi}$ --- which cost the finer
telescoping of \cref{thm:b2:comparison:stirling}. Knowing
which precision each tool buys is half of asymptotic craft.
\end{example}

\begin{example}[Cancellation demands expansions]\label{ex:b2:comparison:cancellation}
Compute the limit of $\sqrt{n^2 + n} - n$. Both terms are
$\sim n$, and ``$\sim n - n$'' is meaningless: equivalents
cannot be subtracted. Expand instead:
\[
\sqrt{n^2 + n} - n
= n\Bigl(\sqrt{1 + \tfrac1n} - 1\Bigr)
= n\Bigl(\frac{1}{2n} - \frac{1}{8n^2} +
O\Bigl(\frac1{n^3}\Bigr)\Bigr)
= \frac12 - \frac{1}{8n} + O\Bigl(\frac1{n^2}\Bigr) :
\]
limit $\frac12$, with the approach speed $\frac1{8n}$ as a
bonus. The mechanism deserves a name: a difference of two large
equivalent quantities lives entirely in their \emph{next}
terms, so one must expand to the first order at which the two
sides differ --- and carry the remainder to certify that
nothing else survives at that order.
\end{example}

\begin{example}[A divergent comparison, worked]\label{ex:b2:comparison:loglog}
For $f(t) = \frac{1}{t\ln t}$ on $\intco{2}{+\infty}$
(\omterm{def:b2:metric:continuity}{continuous}, positive, decreasing): $\int_2^x f = \ln\ln x -
\ln\ln 2 \to \infty$, so by
\cref{thm:b2:comparison:seriesintegral}~(2),
\[
\sum_{k=2}^{n}\frac{1}{k\ln k} = \ln\ln n + C + o(1)
\]
for some constant $C$. Two lessons. First, the divergence is
real but glacial: the partial sum first exceeds $4$ around $n
\approx \eu^{\eu^{4 - C}}$, astronomically large. Second, the
\emph{shape} $\ln\ln n$ was delivered by an antiderivative, not
guessed: for monotone terms, the integral is the canonical
summing device, and the constant $C$ --- like Euler's $\gamma$
--- is the memory of the initial terms.
\end{example}

\section{Stirling's formula}

\begin{lemma}[Wallis integrals, revisited]\label{lem:b2:comparison:wallis}
Let $W_n = \int_0^{\pi/2} \sin^n t\,\dd t$. Then $nW_nW_{n-1} =
\frac\pi2$ for $n \geq 1$, $(W_n)$ decreases, and $W_n \sim
\sqrt{\dfrac{\pi}{2n}}$.
\end{lemma}

\begin{proof}
Integration by parts gives $nW_n = (n-1)W_{n-2}$ ($n \geq 2$), so
$nW_nW_{n-1}$ is constant in $n$, equal to $1 \cdot W_1 W_0 =
\frac\pi2$. Decrease: $\sin^{n+1} \leq \sin^n$ on
$\intcc{0}{\frac\pi2}$. The squeeze, in detail: monotonicity
gives $W_{n+1} \leq W_n \leq W_{n-1}$, and dividing by $W_{n-1}
> 0$,
\[
\frac{n}{n+1} = \frac{W_{n+1}}{W_{n-1}} \leq
\frac{W_n}{W_{n-1}} \leq 1 ,
\]
the left identity from the recurrence at index $n + 1$. Both
bounds tend to $1$: $W_n \sim W_{n-1}$, whence
\[
nW_n^2 \sim nW_nW_{n-1} = \frac\pi2
\qquad\Longrightarrow\qquad
W_n \sim \sqrt{\frac{\pi}{2n}} .
\]
\end{proof}

\begin{example}[The first Wallis integrals]\label{ex:b2:comparison:wallisvalues}
From $W_0 = \frac\pi2$, $W_1 = 1$ and the recurrence $nW_n =
(n-1)W_{n-2}$:
\[
W_2 = \frac\pi4, \qquad
W_3 = \frac23, \qquad
W_4 = \frac{3\pi}{16}, \qquad
W_5 = \frac{8}{15}, \qquad
W_6 = \frac{5\pi}{32}.
\]
Even indices carry a $\pi$, odd ones are rational --- the two
interleaved products of the closed forms. Numerically $W_6
\approx 0.4909$ against the asymptotic
$\sqrt{\pi/12} \approx 0.5116$: at $n = 6$ the equivalent is
already within $5\%$, and the product identity is exact at every
$n$: $6\,W_6W_5 = 6\cdot\frac{5\pi}{32}\cdot\frac8{15} =
\frac\pi2$. Small tables like this one are the cheapest way to
catch an algebra slip before it infects an asymptotic argument.
\end{example}

\begin{theorem}[Stirling]\label{thm:b2:comparison:stirling}
\[
n! \;\sim\; \sqrt{2\pi n}\, \Bigl(\frac{n}{\eu}\Bigr)^{\!n}
.\index{Stirling's formula}
\]
\end{theorem}

\begin{proof}
\emph{Step 1: $n! \sim C \sqrt n\, (n/\eu)^n$ for some constant $C >
0$.} Set
\[
d_n = \ln(n!) - \Bigl(n + \frac12\Bigr)\ln n + n .
\]
Then
\[
d_n - d_{n+1}
= \Bigl(n + \frac12\Bigr) \ln\frac{n+1}{n} - 1
= \Bigl(n + \frac12\Bigr)\Bigl(\frac1n - \frac{1}{2n^2} +
\frac{1}{3n^3} + o\bigl(n^{-3}\bigr)\Bigr) - 1
= \frac{1}{12n^2} + o\Bigl(\frac{1}{n^2}\Bigr),
\]
by the Taylor expansion of $\ln(1 + \frac1n)$. The series $\sum (d_n
- d_{n+1})$ thus converges absolutely (comparison with $\sum
n^{-2}$), so $(d_n)$ converges, say to $d$; exponentiating, $n! \sim
C\sqrt n\,(n/\eu)^n$ with $C = \eu^{d}$.

\emph{Step 2: $C = \sqrt{2\pi}$ via Wallis.} The closed form $W_{2p}
= \frac{(2p)!}{4^p (p!)^2}\cdot\frac\pi2$ (from the recurrence,
Year 1 computation redone in \cref{lem:b2:comparison:wallis}'s
setting) combines with Step 1:
\[
W_{2p} \sim \frac{C\sqrt{2p}\,(2p/\eu)^{2p}}
{4^p\,\bigl(C\sqrt p\,(p/\eu)^p\bigr)^2}\cdot\frac{\pi}{2}
= \frac{\sqrt{2p}}{C\,p}\cdot\frac{\pi}{2}
= \frac{\pi}{C}\cdot\frac{1}{\sqrt{2p}} .
\]
Comparing with $W_{2p} \sim \sqrt{\frac{\pi}{4p}}$
(\cref{lem:b2:comparison:wallis}): $\frac{\pi}{C\sqrt{2p}} =
\sqrt{\frac{\pi}{4p}}\,(1 + o(1))$ forces $C = \pi
\sqrt{\frac{4p}{2p\,\pi}} = \sqrt{2\pi}$.
\end{proof}

\begin{example}[Central binomial coefficient]\label{ex:b2:comparison:centralbinomial}
\[
\binom{2n}{n} = \frac{(2n)!}{(n!)^2}
\sim \frac{\sqrt{4\pi n}\,(2n/\eu)^{2n}}{2\pi n\,(n/\eu)^{2n}}
= \frac{4^n}{\sqrt{\pi n}} :
\]
the probability that a symmetric random walk returns to $0$ at time
$2n$ is $\sim \frac{1}{\sqrt{\pi n}}$ --- an announcement of
\cref{ch:b2:randomvar}.
\end{example}

\begin{remark}[Perspectives within this volume]\label{rem:b2:comparison:perspectives}
Every quantitative chapter ahead speaks this chapter's
language. \cref{ch:b2:series} classifies series by comparing
terms to the scale $n^{-\alpha}(\ln n)^{-\beta}$ --- its
weekend problem maps that frontier completely.
\cref{ch:b2:integration} does the same for improper integrals,
with the identical scale in the \omterm{def:b2:metric:continuity}{continuous} variable.
\cref{ch:b2:powerseries} computes radii of convergence from
$\limsup\abs{a_n}^{1/n}$, an equivalent-of-$n$-th-roots
exercise where Stirling is the standard key ($\sqrt[n]{n!}
\sim \frac n\eu$, \cref{exo:b2:comparison:4}). And the
probability chapters cash Stirling directly: the local
estimates of \cref{ch:b2:randomvar} for binomial coefficients
are \cref{ex:b2:comparison:centralbinomial} and
\cref{ex:b2:comparison:lopsided} verbatim. Asymptotics is not
a chapter here; it is the volume's accent.
\end{remark}

\begin{method}[The bootstrap checklist]\label{met:b2:comparison:checklist}
Before trusting a bootstrapped expansion, audit four points.
(1) \emph{Existence first:} the root or sequence must be pinned
down (monotonicity, intermediate values) before any expansion
--- symbols without referents expand beautifully and mean
nothing. (2) \emph{One order per pass:} each substitution may
only be trusted to the order of the estimate fed in; extracting
two new terms from one pass is the classic source of wrong
coefficients. (3) \emph{Remainders ride along:} carry the
$o(\cdot)$ through every algebraic step and let absorption
(smaller terms swallowed by larger remainders) happen at the
end, explicitly. (4) \emph{Numerical audit:} evaluate at one
honest value of $n$; a coefficient error survives algebraic
re-derivation surprisingly often, and almost never survives
arithmetic.
\end{method}

\begin{remark}[Common pitfalls]\label{rem:b2:comparison:pitfalls}
(i) Equivalents \emph{add} badly: from $u_n \sim n + \ln n$ and
$v_n \sim -n$ one may \emph{not} conclude $u_n + v_n \sim \ln
n$; cancellations demand expansions with explicit remainders,
never bare equivalents. (ii) Never exponentiate an equivalence:
$n + 1 \sim n$ but $\eu^{n+1} \not\sim \eu^n$; the safe
direction is taking logarithms of equivalents tending to
$+\infty$ (this chapter's weekend problem, question 24). (iii)
An \omterm{def:b2:comparison:expansion}{asymptotic expansion} is attached to a \emph{scale}: writing
$f = \frac1x + o\bigl(\frac1{x^2}\bigr)$ claims more than $f =
\frac1x + o\bigl(\frac1x\bigr)$, and mixing the two invalidates
subsequent algebra. (iv) In bootstraps, substitute the
\emph{whole} current expansion, remainder included --- dropping
an $o(\cdot)$ mid-pass produces plausible but wrong
coefficients. (v) The series--integral comparison needs
monotonicity: for oscillating terms it fails outright (compare
$\sum\frac{\sin k}k$, \cref{ch:b2:series}).
\end{remark}

\begin{example}[Stirling in numbers]\label{ex:b2:comparison:stirlingnum}
At $n = 10$: the formula gives $\sqrt{20\pi}\,(10/\eu)^{10}
\approx 3\,598\,696$ against $10! = 3\,628\,800$: relative error
$8.3\cdot10^{-3}$, remarkable for an ``asymptotic'' statement at
$n = 10$. The error has a structure --- the exact refinement $n!
= \sqrt{2\pi n}\,(n/\eu)^n\bigl(1 + \frac1{12n} +
O(n^{-2})\bigr)$ --- whose first correction $\frac1{120} \approx
8.3\cdot10^{-3}$ explains the observed gap almost exactly. The
weekend problem's Euler--Maclaurin machinery is precisely the
systematic source of such correction terms.
\end{example}

\begin{remark}[Where this chapter is used]
Asymptotic comparison is the grammar of everything quantitative
downstream: the convergence tests and Bertrand panorama of
\cref{ch:b2:series}, the integrability criteria of
\cref{ch:b2:integration}, the radius-of-convergence computations
of \cref{ch:b2:powerseries}, and the limit theorems of
\cref{ch:b2:randomvar} (where Stirling runs the de
Moivre--Laplace estimates). The Year 3 volume industrializes the
one idea we prove by hand here --- extract the main term, bound
the rest --- into the Laplace method and dominated convergence.
\end{remark}

\begin{example}[An integral compared to itself: $\int_2^x \frac{\dd t}{\ln t}$]\label{ex:b2:comparison:liworked}
The comparison toolbox also runs on integrals. Let $F(x) =
\int_2^x\frac{\dd t}{\ln t}$ (the integrand is \omterm{def:b2:metric:continuity}{continuous} on
$\intco2\infty$). Integrate by parts:
\[
F(x) = \Bigl[\frac{t}{\ln t}\Bigr]_2^x +
\int_2^x\frac{\dd t}{(\ln t)^2}
= \frac{x}{\ln x} + O\Bigl(\int_2^x\frac{\dd t}{(\ln
t)^2}\Bigr) + O(1),
\]
and the remainder integral is $o\bigl(\frac{x}{\ln x}\bigr)$:
split it at $\sqrt x$, bounding by
\[
\int_2^{\sqrt x}\frac{\dd t}{(\ln t)^2} \leq \sqrt x
\qquad\text{and}\qquad
\int_{\sqrt x}^{x}\frac{\dd t}{(\ln t)^2} \leq
\frac{x}{(\ln\sqrt x)^2} = \frac{4x}{(\ln x)^2} .
\]
Hence $F(x) \sim \frac{x}{\ln x}$. Readers who met the prime
number theorem in this chapter's weekend problem will recognize
$F$: it is the logarithmic integral, the better estimator of
$\pi(x)$, and the computation shows it agrees with
$\frac{x}{\ln x}$ to first order.
\end{example}

\begin{example}[Stirling on a lopsided binomial]\label{ex:b2:comparison:lopsided}
The same three-factorial routine as for
\cref{ex:b2:comparison:centralbinomial} gives, for
$\binom{3n}{n} = \frac{(3n)!}{n!\,(2n)!}$:
\[
\binom{3n}{n} \sim
\frac{\sqrt{6\pi n}\,(3n/\eu)^{3n}}
{\sqrt{2\pi n}\,(n/\eu)^{n}\cdot\sqrt{4\pi n}\,(2n/\eu)^{2n}}
= \sqrt{\frac{3}{4\pi n}}\,
\Bigl(\frac{27}{4}\Bigr)^{\!n} .
\]
The exponential rate $\frac{27}4 = \frac{3^3}{2^2}$ is
$\eu^{3n\,H(1/3)}$ in the entropy notation of information
theory: lopsided binomials grow strictly slower than the central
$4^n$ per two steps --- here $(27/4)^{1/3} \approx 1.89 < 2$ per
step. Every binomial asymptotic in combinatorics and probability
(\cref{ch:b2:randomvar}) is this one computation with different
weights.
\end{example}

\section{Implicitly defined sequences}

\begin{method}\label{met:b2:comparison:implicit}
To find the asymptotics of solutions $x_n$ of an equation $F(x, n) =
0$:\index{implicit sequence}
\begin{enumerate}
  \item \emph{Localize}: prove existence and uniqueness of $x_n$ in
        a definite interval (monotonicity, intermediate value
        theorem), and find its crude behavior (limit, order of
        growth).
  \item \emph{Bootstrap}: substitute the crude form $x_n = (\text{main
        term})(1 + \varepsilon_n)$ into the equation and solve for
        the next order of $\varepsilon_n$; repeat, each pass
        refining one order.
\end{enumerate}
\end{method}

\begin{example}\label{ex:b2:comparison:tan}
For $n \geq 1$, the equation $\tan x = x$ has exactly one solution
$x_n$ in $\intoo{n\pi - \frac\pi2}{n\pi + \frac\pi2}$ (the function
$\tan x - x$ increases from $-\infty$ to $+\infty$ there, its
derivative being $\tan^2 x \geq 0$). \emph{Crude:} $x_n = n\pi +
\frac\pi2 - y_n$ with $y_n \in \intoo{0}{\pi}$; since $x_n \to
\infty$ and $\tan x_n = x_n \to +\infty$, $x_n$ approaches the
asymptote from the left: $y_n \to 0$. \emph{Bootstrap:} $\tan x_n =
\cot y_n = \frac{1}{\tan y_n} \sim \frac{1}{y_n}$, and the equation
$\cot y_n = x_n \sim n\pi$ gives $y_n \sim \frac{1}{n\pi}$. Hence
\[
x_n = n\pi + \frac\pi2 - \frac{1}{n\pi} +
o\Bigl(\frac1n\Bigr),
\]
and the process continues to any order
(\cref{exo:b2:comparison:6}).
\end{example}

\begin{example}[A second run of the method]\label{ex:b2:comparison:xplusln}
Solve $x + \ln x = n$ asymptotically. \emph{Localize:} $x
\mapsto x + \ln x$ increases from $-\infty$ to $+\infty$ on
$\intoo{0}{+\infty}$: a unique root $x_n$, and $x_n \to \infty$.
\emph{Crude:} $\ln x_n = o(x_n)$ gives $x_n \sim n$.
\emph{Bootstrap:} from $x_n = n - \ln x_n$ and $\ln x_n = \ln n
+ o(1)$ (logarithms of equivalents, both sides $\to \infty$):
\[
x_n = n - \ln n + o(1) ;
\]
one more pass, with $\ln x_n = \ln\bigl(n - \ln n + o(1)\bigr) =
\ln n - \frac{\ln n}{n} + o\bigl(\frac{\ln n}n\bigr)$:
\[
x_n = n - \ln n + \frac{\ln n}{n} +
o\Bigl(\frac{\ln n}{n}\Bigr).
\]
(Check at $n = 100$: the root is $x \approx 95.4415$; the
three-term formula gives $100 - 4.6052 + 0.0461 = 95.4409$, the
two-term one $95.3948$ --- each pass gains the predicted
order.) Same loop, third landscape: the method of
\cref{met:b2:comparison:implicit} does not care what the
equation looks like, only that each pass isolates the dominant
unknown.
\end{example}

\section{Exercises}

\begin{exercise}[$\star$]\label{exo:b2:comparison:1}
Expand at $+\infty$, two terms beyond the leading one:
\[
\sqrt{x^2 + x + 1} ,
\qquad
\ln(x^2 + x) - 2\ln x,
\qquad
\frac{x + \sin x}{x - \ln x} .
\]
\end{exercise}

\begin{exercise}[$\star$]\label{exo:b2:comparison:2}
Give the nature (convergence/divergence) and, when divergent, the
leading asymptotics of $\sum_{k \leq n} k^\alpha$ for $\alpha >
-1$, $\alpha = -1$, $\alpha < -1$, via
\cref{thm:b2:comparison:seriesintegral}.
\end{exercise}

\begin{exercise}[$\star\star$]\label{exo:b2:comparison:3}
Prove $H_n = \ln n + \gamma + \frac{1}{2n} + o\bigl(\frac1n\bigr)$.
\emph{(Study $v_n = H_n - \ln n - \gamma$: show $v_n - v_{n+1} =
\frac{1}{2n^2} + O(n^{-3})$ and sum the tail, comparing with
$\sum_{k \geq n} \frac{1}{2k^2} \sim \frac{1}{2n}$ ---
\cref{thm:b2:comparison:seriesintegral}~(1).)}
\end{exercise}

\begin{exercise}[$\star\star$]\label{exo:b2:comparison:4}
Using Stirling, find equivalents of: $\dfrac{(3n)!}{(n!)^3}$;
$\;\dfrac{n!}{n^n}$; $\;\sqrt[n]{n!}$ (as $\frac n\eu(1 + o(1))$,
made precise to two terms).
\end{exercise}

\begin{exercise}[$\star\star$]\label{exo:b2:comparison:5}
For $n \geq 2$, prove that $x^n + x = 1$ has a unique solution $x_n
\in \intoo{0}{1}$, that $x_n \to 1$, and establish
\[
x_n = 1 - \frac{\ln n}{n} + o\Bigl(\frac{\ln n}{n}\Bigr).
\]
\emph{(From $x_n^n = 1 - x_n$: take logarithms and bootstrap with
$x_n = 1 - \varepsilon_n$.)}
\end{exercise}

\begin{exercise}[$\star\star$]\label{exo:b2:comparison:6}
Push \cref{ex:b2:comparison:tan} one order further:
\[
x_n = n\pi + \frac\pi2 - \frac{1}{n\pi} + \frac{1}{2n^2\pi} +
o\Bigl(\frac{1}{n^2}\Bigr).
\]
\emph{(Write $\cot y_n = x_n$ exactly, expand $\cot y = \frac1y -
\frac y3 + o(y)$ and $x_n = n\pi(1 + \frac{1}{2n} - \dots)$, and
identify.)}
\end{exercise}

\begin{exercise}[$\star\star$]\label{exo:b2:comparison:7}
Determine $\lim_{n\to\infty} \dfrac{1}{n!}\sum_{k=0}^{n} k!$
\emph{(bound the sum of all terms but the last two)}, and deduce the
\omterm{def:b2:comparison:expansion}{asymptotic expansion} $\sum_{k \leq n} k! = n!\bigl(1 + \frac1n +
O(n^{-2})\bigr)$.
\end{exercise}

\begin{exercise}[$\star\star\star$]\label{exo:b2:comparison:8}
Let $u_0 > 0$ and $u_{n+1} = u_n + \dfrac{1}{u_n}$. Prove that $u_n
\to \infty$, then that $u_n \sim \sqrt{2n}$ \emph{(study $u_n^2$:
its increments are $2 + u_n^{-2}$; sum)}, and refine:
\[
u_n = \sqrt{2n}\Bigl(1 + \frac{\ln n}{8n} +
o\Bigl(\frac{\ln n}{n}\Bigr)\Bigr).
\]
\emph{(From $u_n^2 = 2n + \sum_{k<n} u_k^{-2} + u_0^2$ and $u_k^2
\sim 2k$: the sum is $\sim \frac12\ln n$ by
\cref{thm:b2:comparison:seriesintegral}.)}
\end{exercise}

\begin{exercise}[$\star\star\star$]\label{exo:b2:comparison:9}
(A Riemann sum with a twist) Determine the asymptotic behavior of
\[
S_n = \sum_{k=1}^{n} \frac{1}{n + k\ln n} .
\]
\emph{(Factor $n$: $S_n = \frac1n\sum_k \bigl(1 +
\frac{k\ln n}{n}\bigr)^{-1}$; recognize a Riemann-type sum with a
slowly varying parameter $t = \ln n$, compute $\int_0^1
\frac{\dd u}{1 + tu} = \frac{\ln(1+t)}{t}$, and conclude $S_n \sim
\frac{\ln\ln n}{\ln n}$.)}
\end{exercise}

\begin{exercise}[$\star$]\label{exo:b2:comparison:10}
Prove the identity $(\ln n)^{\ln n} = n^{\ln\ln n}$, then rank the
following in increasing $o(\cdot)$ order at infinity, with
proofs: $n^2$, $(\ln n)^{\ln n}$, $2^n$, $n!$, $n^n$.
\end{exercise}

\begin{exercise}[$\star\star$]\label{exo:b2:comparison:11}
(Tail of $\sum 1/k^2$, two terms) Using the exact telescoping
$\sum_{k > n} \frac{1}{k(k+1)} = \frac{1}{n+1}$ and the
decomposition $\frac1{k^2} = \frac{1}{k(k+1)} +
\frac{1}{k^2(k+1)}$, prove
\[
\sum_{k > n} \frac{1}{k^2}
= \frac1n - \frac{1}{2n^2} + O\Bigl(\frac{1}{n^3}\Bigr).
\]
\end{exercise}

\begin{exercise}[$\star\star\star$]\label{exo:b2:comparison:12}
Let $u_0 = \frac12$ and $u_{n+1} = u_n + \eu^{-u_n}$. Prove that
$u_n \to \infty$, then --- setting $v_n = \eu^{u_n}$ and showing
$v_{n+1} = v_n + 1 + \frac{1}{2v_n} + O\bigl(v_n^{-2}\bigr)$ ---
establish
\[
u_n = \ln n + \frac{\ln n}{2n} + O\Bigl(\frac1n\Bigr).
\]
\end{exercise}

\section{Problem: Bootstrapping, from Euler--Maclaurin to the
Primes}

An implicit or accumulated quantity rarely hands over its
asymptotics at once; one extracts them in passes, each pass
feeding the previous estimate back into the defining relation.
This weekend problem trains that loop on fresh equations, proves
the first-order \emph{Euler--Maclaurin formula} (the trapezoid
upgrade of the series--integral comparison, with rigorous error
bars), inverts $x\ln x = n$, and cashes the method's most famous
cheque: from the admitted prime number theorem, the asymptotic
law $p_n \sim n\ln n$ of the $n$-th prime.

\begin{problem}[{Weekend problem --- the Euler--Maclaurin
correction and the asymptotics of the $n$-th prime}]
\label{pb:b2:comparison:1}

\medskip
\noindent\textbf{Part I --- The bootstrap loop on a fresh
equation.}
\begin{enumerate}
  \item Prove the uniqueness claim of
        \cref{def:b2:comparison:expansion}: if $f = \sum_{i\leq
        k} c_i\varphi_i + o(\varphi_k) = \sum_{i \leq k}
        c_i'\varphi_i + o(\varphi_k)$ along the same scale, then
        $c_i = c_i'$ for all $i$. Then push the course's mixed
        example one rung further:
        \[
        \frac{1}{x - \ln x} = \frac1x + \frac{\ln x}{x^2} +
        \frac{(\ln x)^2}{x^3} + o\Bigl(\frac{(\ln
        x)^2}{x^3}\Bigr) \qquad (x \to +\infty),
        \]
        and explain why no term $\frac{c}{x^2}$ appears.
  \item Show that for every $n \geq 1$ the equation $\eu^x + x =
        n$ has exactly one real solution $x_n$, and that $x_n
        \to +\infty$ with $x_n \sim \ln n$.
  \item Bootstrap twice:
        \[
        x_n = \ln n - \frac{\ln n}{n} - \frac{(\ln n)^2}{2n^2} +
        o\Bigl(\frac{(\ln n)^2}{n^2}\Bigr).
        \]
  \item Check numerically at $n = 1000$: compare $x_{1000}
        \approx 6.90083$ with the one-, two- and three-term
        values of question 3, to five decimals.
\end{enumerate}

\medskip
\noindent\textbf{Part II --- Euler--Maclaurin, order one.}
\begin{enumerate}[resume]
  \item Prove the trapezoid kernel identity: for $g$ of class
        $C^2$ on $\intcc{0}{1}$,
        \[
        \int_0^1 g(t)\,\dd t = \frac{g(0) + g(1)}{2}
        - \frac12\int_0^1 t(1 - t)\,g''(t)\,\dd t
        \]
        \emph{(integrate $\frac12 t(1-t)g''$ by parts twice)}.
  \item Let $f$ be $C^2$ on $\intco{1}{+\infty}$ with
        $\int_1^\infty \abs{f''} < \infty$. Show that
        \[
        E_n = \sum_{k=1}^{n} f(k) - \int_1^n f -
        \frac{f(1) + f(n)}{2}
        \]
        converges to a constant $E$, with the tail bound
        $\abs{E - E_n} \leq \frac18\int_n^\infty\abs{f''}$: the
        \emph{Euler--Maclaurin formula} to first order.
  \item Apply this to $f(t) = \frac1t$: prove
        \[
        H_n = \ln n + \gamma + \frac{1}{2n} + \varepsilon_n,
        \qquad \abs{\varepsilon_n} \leq \frac{1}{8n^2},
        \]
        strengthening \cref{exo:b2:comparison:3} (identify the
        constant with $\gamma$ by comparing with
        \cref{ex:b2:comparison:harmonic}).
  \item Extract the next coefficient: show $\varepsilon_n =
        -\frac{1}{12n^2} + o\bigl(\frac1{n^2}\bigr)$
        \emph{(the increments of $E_n$ are
        $\frac12\int_0^1t(1-t)f''(n+t)\dd t =
        \frac1{12}f''(n) + o(f''(n))$; sum the tail with
        \cref{thm:b2:comparison:seriesintegral})}.
  \item Apply question 6 to $f = \ln$: re-derive in three lines
        the convergence of $d_n = \ln n! - (n +
        \frac12)\ln n + n$ (Step 1 of
        \cref{thm:b2:comparison:stirling}), with the bonus error
        rate $d_n = d + O\bigl(\frac1n\bigr)$.
  \item Apply question 6 to $f(t) = \frac{1}{\sqrt t}$: show
        \[
        \sum_{k=1}^{n}\frac1{\sqrt k} = 2\sqrt n + c +
        \frac{1}{2\sqrt n} + O\Bigl(\frac{1}{n^{3/2}}\Bigr)
        \]
        for some constant $c$, and evaluate all terms at $n =
        10^4$ (the constant is $c \approx -1.4604$).
\end{enumerate}

\medskip
\noindent\textbf{Part III --- Inversion: the equation $x\ln x =
n$.}
\begin{enumerate}[resume]
  \item Show that $x\ln x = n$ has exactly one solution $x_n \in
        \intco{1}{+\infty}$ for $n \geq 1$, that $x_n \to
        \infty$, and that $\ln x_n \sim \ln n$.
  \item Deduce the one-term inversion $x_n \sim
        \dfrac{n}{\ln n}$, then bootstrap once more:
        \[
        \ln x_n = \ln n - \ln\ln n + o(1),
        \qquad
        x_n = \frac{n}{\ln n}\Bigl(1 + \frac{\ln\ln n}{\ln n}
        + o\Bigl(\frac{\ln\ln n}{\ln n}\Bigr)\Bigr).
        \]
  \item Test at $n = 10^6$: the true root is $x \approx
        87\,848$; compare with the one-term ($\approx 72\,382$)
        and two-term ($\approx 86\,140$) values, and explain the
        slow gain (the expansion parameter is $\frac{\ln\ln
        n}{\ln n}$, only $\approx 0.19$ at $n = 10^6$).
  \item We now \emph{admit} the prime number theorem: the number
        $\pi(x)$ of primes $\leq x$ satisfies $\pi(x) \sim
        \frac{x}{\ln x}$ as $x \to \infty$ (proved honestly in
        the Year 3 volume). Writing $p_n$ for the $n$-th prime,
        justify $\pi(p_n) = n$, and run the inversion of
        questions 11--12 to prove
        \[
        p_n \sim n \ln n .
        \]
  \item Dividends: (a) show $\sum_{k \leq n} p_k \sim
        \frac{n^2\ln n}{2}$ \emph{(compare $\sum k\ln k$ with
        $\int t\ln t\,\dd t$)}; (b) compute the approximate
        chance that a uniformly random integer with $100$
        digits is prime ($\ln 10^{100} \approx 230.26$: about
        one in $230$).
\end{enumerate}

\medskip
\noindent\textbf{Part IV --- The method exported: $x\tan x =
1$.}
\begin{enumerate}[resume]
  \item Show that for each $n \geq 1$ the equation $\tan x =
        \frac1x$ has exactly one solution $x_n$ in
        $\intoo{n\pi}{\,n\pi + \frac\pi2}$, and that $z_n = x_n
        - n\pi \to 0^+$.
  \item One term: $z_n \sim \dfrac{1}{n\pi}$.
  \item Show that the expansion of $z_n$ has \emph{no}
        $\frac{c}{n^2}$ term: $z_n = \frac1{n\pi} +
        O\bigl(\frac{1}{n^3}\bigr)$.
  \item Three terms: using $\arctan u = u - \frac{u^3}3 +
        O(u^5)$ and $\frac1{x_n} = \frac{1}{n\pi} -
        \frac{z_n}{(n\pi)^2} + O(n^{-3}\cdot z_n^2)$, prove
        \[
        x_n = n\pi + \frac{1}{n\pi} -
        \frac{4}{3\pi^3 n^3} + o\Bigl(\frac{1}{n^3}\Bigr).
        \]
  \item Check at $n = 3$: true root $x_3 \approx 9.5293344$;
        compare the one- and three-term values, and contrast in
        one sentence with the course's $\tan x = x$
        (\cref{ex:b2:comparison:tan}): where each sequence sits
        in its window, and why.
\end{enumerate}

\medskip
\noindent\textbf{Part V --- A dynamical bootstrap, rules of the
game, synthesis.}
\begin{enumerate}[resume]
  \item Let $u_0 \in \intoo{0}{\pi}$ and $u_{n+1} = \sin u_n$.
        Show $u_n \to 0$ decreasingly, and compute the limit of
        $\dfrac{1}{u_{n+1}^2} - \dfrac{1}{u_n^2}$ \emph{(expand
        $\sin^{-2}$ via $\sin u = u - \frac{u^3}6 + o(u^3)$)}.
  \item Deduce, via Ces\`aro means (Year 1 volume), the classic
        \[
        u_n \sim \sqrt{\frac{3}{n}} .
        \]
  \item (Certified numerics) Using question 7's rigorous bound,
        show that evaluating $\ln n + \gamma + \frac1{2n}$ at $n
        = 10^6$ yields $H_{10^6}$ with error at most
        $1.25\cdot10^{-13}$ --- a million-term sum computed to
        thirteen digits by three terms.
  \item (Rules of the game) Prove or refute, with proofs or
        counterexamples: (a) if $u_n \sim v_n \to +\infty$ then
        $\ln u_n \sim \ln v_n$; (b) if $u_n \sim v_n$ then
        $\eu^{u_n} \sim \eu^{v_n}$; (c) if $f \sim g$ at
        $+\infty$ ($f, g$ differentiable) then $f' \sim g'$.
  \item (Synthesis) In one sentence each: the bootstrap loop of
        \cref{met:b2:comparison:implicit} as used in Parts I,
        III, IV; what the trapezoid correction adds to
        \cref{thm:b2:comparison:seriesintegral}; why inversion
        of $x\ln x$ is exactly the bridge from $\pi(x)$ to
        $p_n$; and which of question 24's rules protected which
        step. Name the two summits: the Euler--Maclaurin formula
        (first order), and the asymptotic law of the $n$-th
        prime.
\end{enumerate}
\end{problem}
